\chapter{Aortic Stenosis}

\section*{Definition and Overview}
Aortic stenosis is a heart valve condition in which the aortic valve -- the gateway between the heart's main pumping chamber (the left ventricle) and the aorta -- becomes narrowed, thickened, and stiff, so the heart must work much harder to push blood out into the body. Over time this extra workload thickens the heart muscle (left ventricular hypertrophy) and eventually leads to heart failure, angina, and fainting.

It is the most common valve disease in older adults. Many people are symptom-free for years while the valve narrows gradually, and the condition is often first detected as a heart murmur. Once symptoms develop, however, the outlook without treatment is poor, and valve replacement is usually needed. The two modern approaches to replacement are surgical valve replacement (SAVR) and transcatheter aortic valve implantation (TAVI), a less invasive procedure.

\section*{Epidemiology}
Aortic stenosis affects roughly one in twenty adults over the age of 75 and is the most common heart valve problem in developed countries. Milder thickening of the valve without significant obstruction -- aortic sclerosis -- is even more common, affecting up to a quarter of older adults. The condition is slightly more common in men than in women.

The dominant cause, age-related calcification, means prevalence rises steeply with age. In younger people, the usual cause is a congenital bicuspid aortic valve, which calcifies and narrows decades earlier than a normal three-leaflet valve. In parts of the world where rheumatic fever remains common, rheumatic heart disease is a leading cause of aortic stenosis in younger adults.

\section*{Aetiology and Pathophysiology}
The normal aortic valve has three thin leaflets that open fully during systole and close tightly during diastole. Aortic stenosis develops when these leaflets thicken, stiffen, and fail to open completely:

\begin{itemize}
    \item \textbf{Age-related calcification:} calcium deposits progressively build up on the valve over decades, the commonest cause in older adults
    \item \textbf{Bicuspid aortic valve:} a valve with only two leaflets instead of three, present from birth, which tends to calcify and narrow in mid-life
    \item \textbf{Rheumatic fever:} previous rheumatic fever can scar and fuse the valve leaflets, causing combined stenosis and regurgitation
    \item \textbf{Chronic kidney disease:} altered calcium and phosphate handling accelerates valve calcification
    \item \textbf{Risk factors:} high blood pressure, high cholesterol, diabetes, and smoking are associated with faster progression of valve disease
    \item \textbf{Mechanism of symptoms:} the narrowed orifice forces the left ventricle to generate very high pressures, thickening its wall; when the supply of blood to the thickened muscle cannot keep up, angina, fainting, and heart failure follow
\end{itemize}

\section*{Clinical Features}
\begin{itemize}
    \item \textbf{No symptoms for many years:} the valve often narrows gradually while the person remains entirely well
    \item \textbf{Exertional chest pain (angina):} the classic triad of severe disease begins here, as the thickened muscle outgrows its blood supply
    \item \textbf{Fainting on exertion:} blacking out during or shortly after activity, caused by a fall in blood pressure when the obstructed valve cannot increase output
    \item \textbf{Breathlessness on activity:} progressing to heart failure symptoms such as fatigue, swollen ankles, and breathlessness at rest
    \item \textbf{Palpitations:} a fluttering or racing heart, especially if atrial fibrillation develops
    \item \textbf{A heart murmur:} a harsh systolic murmur, best heard over the right upper chest, as blood is forced through the narrow valve
    \item \textbf{A weak, delayed pulse:} the pulse at the wrist is typically small and slow-rising
\end{itemize}

\section*{Differential Diagnosis}
\begin{itemize}
    \item \textbf{Aortic sclerosis:} thickening of the valve without significant obstruction, which requires monitoring but not surgery
    \item \textbf{Hypertrophic cardiomyopathy:} causes chest pain, fainting, and a systolic murmur from muscular obstruction rather than valve disease
    \item \textbf{Other valve problems:} aortic regurgitation and mitral valve disease, which can coexist or mimic aortic stenosis
    \item \textbf{Coronary artery disease:} causes the same symptoms of angina, breathlessness, and heart failure
    \item \textbf{High blood pressure and heart failure from other causes}
\end{itemize}

\section*{When to Seek Medical Care}
See a doctor if you develop chest pain, fainting, breathlessness on exertion, or a racing heartbeat -- these can be the first signs of significant aortic stenosis. Anyone with a known valve problem, a bicuspid valve, or a family history of valve disease should attend regular reviews so the valve can be monitored. Valve disease during pregnancy needs specialist obstetric-cardiology care.

\textbf{Call 000 (or local emergency services) for:}
\begin{itemize}
    \item Severe or persistent chest pain
    \item Fainting, collapse, or near-syncope
    \item Sudden breathlessness, especially at rest
    \item A fast, irregular heartbeat that does not settle
\end{itemize}

\section*{Diagnosis and Investigations}
\begin{itemize}
    \item \textbf{Clinical examination:} a stethoscope detects the characteristic murmur, and the carotid pulse may be weak and delayed
    \item \textbf{Echocardiography:} the key test, measuring the valve area and pressure gradient, assessing left ventricular function, and detecting valve thickening and calcification
    \item \textbf{ECG:} shows left ventricular hypertrophy and any rhythm disturbance
    \item \textbf{Exercise testing:} occasionally performed under supervision to assess symptom response and risk in asymptomatic patients
    \item \textbf{Cardiac CT:} quantifies valve calcium and provides anatomical detail used to plan TAVI
    \item \textbf{Cardiac catheterisation:} sometimes used to measure pressures directly, especially if non-invasive tests are unclear
    \item \textbf{Blood tests:} kidney function, and screening for conditions relevant to surgery
    \item \textbf{Strain and stress imaging:} exercise echocardiography in selected cases to unmask symptoms and define severity
\end{itemize}

\section*{Management}
\begin{itemize}
    \item \textbf{Monitoring:} mild to moderate stenosis is reviewed regularly with echocardiography, and exercise is generally encouraged
    \item \textbf{Medicine:} no drug reverses the narrowing, but medicines treat heart failure symptoms, control blood pressure, and manage rhythm problems
    \item \textbf{Aortic valve replacement:} surgical replacement (SAVR) or transcatheter implantation (TAVI) is the definitive treatment for severe, symptomatic stenosis and is now also offered to selected asymptomatic patients with very severe disease
    \item \textbf{Treatment of risk factors:} control of cholesterol, blood pressure, and diabetes, and smoking cessation, slow the progression of associated vascular disease
    \item \textbf{Dental and infection care:} good dental hygiene and prompt treatment of infections reduce the small risk of valve infection (endocarditis)
    \item \textbf{Antibiotic prophylaxis:} only for specific high-risk groups undergoing certain procedures, as recommended by current guidelines
\end{itemize}

\section*{Complications}
\begin{itemize}
    \item Heart failure, from the long-term pressure load on the left ventricle
    \item Abnormal heart rhythms, especially atrial fibrillation
    \item Infective endocarditis -- infection of the valve, which is serious in a stenotic valve
    \item Stroke from emboli arising on the valve or in the left atrium
    \item Fainting and falls from exertional blackouts
    \item Sudden death, mainly in people with severe symptomatic disease
    \item Complications of valve replacement itself, including bleeding, stroke, pacemaker requirement, and prosthetic valve problems
\end{itemize}

\section*{Prognosis and Outcomes}
People with mild to moderate aortic stenosis can remain stable for many years under surveillance. The outlook changes dramatically once symptoms develop in severe disease: without valve replacement, average survival is typically only a few years, with a high risk of death from heart failure or sudden cardiac death.

The turning point is symptom onset, at which stage valve replacement is strongly recommended. After successful replacement, most people experience marked improvement in symptoms and a substantially better life expectancy. Long-term outcomes depend on age, general health, and the presence of other cardiac disease, but the procedure is life-changing for the majority of patients.

\section*{Prevention and Self-Care}
\begin{itemize}
    \item Maintain a heart-healthy lifestyle: a balanced diet, regular exercise, and a healthy weight
    \item Control blood pressure and cholesterol, and avoid smoking
    \item Attend regular check-ups if you have a bicuspid valve, a family history of valve disease, or a murmur
    \item Keep up good dental hygiene and seek prompt treatment of infections
    \item Report chest pain, fainting, or breathlessness promptly rather than waiting
    \item Tell clinicians and dentists about your valve condition before any planned procedure
    \item Follow specialist advice on exercise intensity and, if relevant, pregnancy planning
\end{itemize}

\section*{Sources and Further Reading}
The following organisations provide reliable, up-to-date information on aortic stenosis for patients, students, and health professionals.

\begin{itemize}
    \item NHS. \textit{Information on aortic stenosis, symptoms, and valve treatment.} https://www.nhs.uk/conditions/
    \item Mayo Clinic. \textit{Overview of aortic stenosis, diagnosis, and treatment.} https://www.mayoclinic.org/
    \item NICE. \textit{Guidance on aortic valve disease and transcatheter aortic valve implantation.} https://www.nice.org.uk/
    \item MedlinePlus. \textit{Health information on aortic stenosis and heart valve disease.} https://medlineplus.gov/
    \item MSD Manual. \textit{Professional overview of aortic stenosis.} https://www.msdmanuals.com/
    \item World Health Organization. \textit{Resources on rheumatic heart disease and cardiovascular health.} https://www.who.int/
\end{itemize}
